% !TeX root = ../main.tex

\chapter{总结与展望}

\section{总结}

本文围绕现实网络中广泛存在的路径依赖现象，研究了路径代价对有限阶历史状态敏感时的最短路径求解问题.针对该类问题在原图上难以直接应用传统加性最短路算法的困难，本文基于线图映射与条件加权思想，构建了条件权重最短路径（Conditional Weighting Shortest Path, CWSP）框架，并进一步从工程实现角度提出了面向 GPU 的三阶段并行化优化方案.全文的主要工作与结论可概括如下.

（1）\textbf{提出并系统阐释了线图映射求解有限阶历史依赖最短路径的统一建模框架.}
本文以“拓扑替代—条件加权—替代图寻优”的流水线形式，将原网络中受历史状态影响的动态代价判定，等价转化为替代图上的静态边权，从而将复杂约束问题规约为标准加性最短路求解.该思路不仅为非加性约束的数学处理提供了结构化路径，也为后续并行化与性能拆解提供了自然边界.

（2）\textbf{实现了可复现的三阶段基准并提出面向 GPU 的关键优化策略.}
针对线图化会引入显著的规模放大与内存压力的问题，本文在三个关键步骤上分别设计并实现了并行优化：阶段1采用前缀和预扫描与无锁 CSR 并发写入，实现大规模图结构的高吞吐构造；阶段2将条件权重计算映射为边级数据并行，并通过分支规整降低线程束内分歧；阶段3采用桶式 $\Delta$-Stepping 算法替代全序优先队列的严格顺序弹出，利用同桶波前并行松弛提升 GPU 侧可并行度.

（3）\textbf{通过多组实验验证了算法正确性与并行方案的阶段性与端到端收益.}
在微观机场路网实验中，线图 CWSP 能稳定规避历史依赖约束导致的不可行解，并与传统算法在可行区间保持一致的代价值表现，验证了建模与求解框架的正确性.在规模化性能实验中，GPU 方案在拓扑构建、条件加权与并行寻优三个阶段均表现出随规模增长而增强的加速趋势，端到端总耗时增长更为平缓，从而证明了三阶段协同优化在大规模替代图上的可扩展性.

（4）\textbf{分析了高阶历史依赖情形下的状态空间膨胀与资源瓶颈.}
当历史阶数 $k$ 提升时，替代图规模 $|V_H|$ 与 $|E_H|$ 会呈指数式增长，阶段1/2 的结构生成与加权代价随之显著抬升；同时，阶段3的寻优也会受到规模放大的系统性影响.实验表明在较高 $k$ 区间 GPU 的端到端加速比进一步放大，但替代图的生成与存储将更早触碰单卡显存上限（例如出现 out-of-memory），提示高阶场景的工程落地需要进一步的结构压缩与资源管理策略.

\section{展望}

尽管本文已在可复现的实验链路上验证了 CWSP 框架与 GPU 并行化方案的有效性，但在更大规模与更高阶历史依赖场景下仍存在进一步改进空间.后续研究可从以下方向展开.

（1）\textbf{面向高阶历史依赖的替代图压缩与显存友好构造.}
高阶 $k$ 下替代图的指数式膨胀会使阶段1的结构生成与存储成为最先触碰硬件上限的瓶颈.未来可研究更紧凑的状态表示、按需展开（lazy expansion）、分层裁剪与外存/流式构造等技术，以在保证等价性的前提下抑制 $|V_H|,|E_H|$ 的爆炸增长.

（2）\textbf{并行寻优算法的进一步增强与鲁棒参数自适应.}
桶式 $\Delta$-Stepping 算法的性能对 $\Delta$ 参数与权重分布具有敏感性.后续可考虑自适应桶宽、混合策略（如在不同尺度区间切换 Dijkstra/$\Delta$-Stepping 算法）、降低原子争用与冗余松弛的调度方法，从而提升在多分布权重与稠密图上的稳定性.

（3）\textbf{面向查询型应用的早停与点对点口径优化.}
本文实验以单源最短路为基准以刻画全图求解能力，但实际场景常为多次点对点查询.未来可在保持可复现计时口径的基础上引入早停策略、双向搜索或分层索引，以更贴近在线调度/路径查询的业务形态，并进一步推动阶段瓶颈从“寻优”向“构图/加权”的前移与优化.

（4）\textbf{多 GPU/分布式扩展与真实场景融合验证.}
面对更大规模网络与更高阶历史依赖需求，可进一步研究多 GPU 下的图划分、跨设备桶同步与通信隐藏策略，并将本文方法与真实交通/航空/机器人路径规划的动态约束模型深度结合，在更复杂的数据与约束组合中验证鲁棒性与工程可用性.

综上，本文工作为路径依赖型最短路径问题提供了从建模到并行求解的可复现实现路径.随着图计算硬件与并行框架的持续演进，面向更高阶历史依赖、更大规模网络与更强工程约束的 CWSP 扩展仍具有广阔的研究与应用空间.
